\section{单位、整除与相伴}


本讲从整数"唯一分解"这件事出发，把它推广到一般整环。推广时有两件事必须先说清楚：\textbf{哪些元素算"可逆"}（单位），以及\textbf{两个元素在什么意义下算"同一个因子"}（相伴）。这一节把这两个概念和整除的关系理清。

\textbf{本节用到的记号，先一次交代清楚：}

\begin{itemize}
\item $R$：整环（交换含幺环且 $1\ne0$、无零因子）；$R^\times$：$R$ 的单位群，即全体可逆元。
\item $a\mid b$：$a$ 整除 $b$，即存在 $c\in R$ 使 $b=ac$。
\item $a\nmid b$：$a$ 不整除 $b$。
\item $(a)=\{ra:r\in R\}$：由 $a$ 生成的主理想。
\item $\mathbb Z[\sqrt{d}]=\{a+b\sqrt{d}:a,b\in\mathbb Z\}$（$d$ 为无平方因子的整数）：$\mathbb Z$ 添上 $\sqrt{d}$ 得到的环，是本讲反复使用的反例来源。
\end{itemize}


\subsection{单位与整除}


含幺环 $R$ 的\textbf{单位群}是

\[
R^\times=\{u\in R:\exists v\in R,\ uv=vu=1\},
\]

即全体有乘法逆元的元素。它在乘法下构成群。$\mathbb Z$ 中 $\mathbb Z^\times=\{\pm1\}$；域 $F$ 中 $F^\times=F\setminus\{0\}$；而 $F[x]$ 中 $(F[x])^\times=F\setminus\{0\}$（非零常数），见 [Lec 09 §4](../lec09/sections/04-units-fractions-characteristic.md)。

\textbf{整除}的定义是

\[
a\mid b\iff \exists c\in R,\ b=ac .
\]

整除有两条常用性质，都由定义直接得到：

\begin{itemize}
\item 传递性：$a\mid b$ 且 $b\mid c$ $\Longrightarrow$ $a\mid c$；
\item 与理想的关系：$a\mid b\iff (b)\subseteq(a)$。\textbf{"整除"是"包含"的反向写法}，这一点在下一节讨论因子链时非常关键。
\end{itemize}


\subsection{相伴：什么时候两个元素算"同一个因子"}


若存在单位 $u\in R^\times$ 使 $a=ub$，就说 $a$ 与 $b$ \textbf{相伴}，记作 $a\approx b$。例如在 $\mathbb Z$ 中 $2\approx-2$，在 $F[x]$ 中 $f\approx cf$（$c\ne0$）。

相伴是一个等价关系（单位群是群，验证自反、对称、传递都用群公理），所以"相伴"把环中的元素分成若干类。为什么要引入它？因为在整环里"$a$ 和 $b$ 互相整除"恰恰等价于"$a$ 和 $b$ 相伴"，而\textbf{不是}等价于 $a=b$：

\textbf{命题。} 设 $R$ 是整环，$a,b\ne0$。则

\[
a\mid b\ \text{且}\ b\mid a\iff a\approx b .
\]

\textbf{证明。} （$\Leftarrow$）若 $a=ub$（$u$ 为单位），则 $b=u^{-1}a$，于是 $a\mid b$ 且 $b\mid a$。

（$\Rightarrow$）设 $b=ac$、$a=bd$（$c,d\in R$）。代入得 $a=acd$，即 $a(1-cd)=0$。因为 $a\ne0$ 且 $R$ 是整环（无零因子），所以 $cd=1$，即 $c,d$ 都是单位。于是 $b=ac$ 且 $c\in R^\times$，故 $a\approx b$。$\square$

\textbf{这条命题为什么重要。} 在 $\mathbb Z$ 里 $6=2\cdot3=(-2)\cdot(-3)$，如果"唯一分解"要求分解式逐字相同，那 6 就有两种分解。但 $2\approx-2$、$3\approx-3$，两式其实只差单位。所以唯一分解的正确说法必须是\textbf{"在相伴意义下唯一"}，这正是下一节要定义的\textbf{唯一分解整环}（Unique Factorization Domain，简称 \textbf{UFD}）里要写进去的。

\subsection{相伴的三种等价刻画}


在整环中以下三件事等价，可以互换使用：

\[
a\approx b\iff a\mid b\ \text{且}\ b\mid a\iff (a)=(b).
\]

第一、二条的等价就是上面的命题；第一条与第三条的等价来自"$a\mid b\iff(b)\subseteq(a)$"：互相整除就是两个主理想互相包含，也就是相等。

\textbf{读题提示：} 看到"$a$ 与 $b$ 相伴"，可以立刻换成"$(a)=(b)$"来算，往往更方便；反过来看到两个主理想相等，也可以换成相伴来谈元素。

\subsection{参考对照}


\begin{itemize}
\item [DF] §8.1（整除、相伴、不可约元与素元的一般讨论）。
\item [DS] §10.1（整环与单位；"相伴"只在 §8.1.2 命题 (vi) 里对 $\mathbb Z$ 定义过）。
\item 课程讲义 §7.3（相伴的定义 7.19）。
\item 原始材料：本地归档（手写原稿与录音转写），不纳入课程仓库。
\end{itemize}
